\clearpage
\section*{Appendix AT. Exact fourfold mediator energy and carrier-coupling criterion}
\addcontentsline{toc}{section}{Appendix AT. Exact fourfold mediator energy and carrier-coupling criterion}

We record the finite Fourier calculation independently of the unresolved
singular-carrier construction. The unit-torus convention is
$e_k(x)=\exp(2\pi i k\cdot x)$, with viscosity $\nu>0$.
The finite triadic transfer formula itself is standard
\cite{DarVermaEswaran2001}; the exact arithmetic below concerns a
particular fourfold family, not a general cascade theorem.

\paragraph{The fourfold real seed.}
Set
\[
 p=(1,0,0),\quad q=(0,2,1),\quad r=p+q=(1,2,1),
 \qquad a=(0,1,0),\quad b=(0,1,-2).
\]
For real amplitudes $A,B,C$ and $\sigma\in\{-1,+1\}$ prescribe
\[
 \widehat h(p)=iAa,\qquad \widehat h(q)=Bb,\qquad
 \widehat h(r)=\sigma Cb.
\]
Extend these coefficients by complex conjugation at $-k$ and by
the four quarter-turns about the $z$ axis, rotating both wavevectors
and polarizations. The resulting real solenoidal trigonometric
polynomial $h$ has exactly 20 distinct possible wavevectors:
four of type $p$ and eight each of types $q,r$.
The $p$ coefficients counted under opposite quarter-turns coincide,
so they are not duplicated in the Fourier series. Let
$\mathcal S=\operatorname{supp}\widehat h$ for generic nonzero amplitudes,
and define
\[
 \mathcal B_k(h,h)=\sum_{\ell+m=k}
  (\widehat h(\ell)\cdot m)\,\widehat h(m),\qquad
 \mathbb P_k=I-\frac{k\otimes k}{|k|^2}
 \quad(k\ne0).
\]

\noindent\textbf{Proposition AT.1 (finite mediator enumeration and leakage).}
Treat $A,B,C$ as independent formal real variables. Among all 170
distinct nonzero pair sums, 124 have a nonzero Leray-projected
quadratic coefficient polynomial. Of these, 108 lie outside $\mathcal S$.
Denote this last set by $\mathcal N$. It decomposes into 92
cross-sector-only modes, which admit no contributing pair in one
unrotated sector, and 16 modes with contributing pairs entirely within
an unrotated sector. Relative to the highest seed wavenumber
$|r|^2=6$, there are 72 new modes with $|k|^2>6$, 8 with
$|k|^2=6$, and 28 with $|k|^2<6$.

The exact off-support quadratic source norm is
\begin{equation}\label{eq:AT-leakage}
\begin{aligned}
 \mathcal L(A,B,C)
 &=\sum_{k\in\mathcal N}|\mathbb P_k\mathcal B_k(h,h)|^2\\
 &=176A^2B^2+\frac{8000}{33}A^2C^2
  +\frac{1600}{3}B^4
  +\frac{1796592}{1105}B^2C^2
  +\frac{25376}{35}C^4 .
\end{aligned}
\end{equation}
For the 92 cross-sector-only modes, the corresponding sum is
\begin{equation}\label{eq:AT-cross}
 \mathcal L_\times
 =16A^2B^2+\frac{2720}{33}A^2C^2
 +\frac{1600}{3}B^4+\frac{1796592}{1105}B^2C^2
 +\frac{25376}{35}C^4,
\end{equation}
while the 16 within-sector modes contribute
$160A^2(B^2+C^2)$. The strictly higher-frequency subfamily has
\begin{equation}\label{eq:AT-high}
 \mathcal L_{>6}
 =8A^2B^2+\frac{2256}{11}A^2C^2
 +\frac{1600}{3}B^4+\frac{262432}{221}B^2C^2
 +\frac{25376}{35}C^4 .
\end{equation}
All three polynomials are independent of the phase sign $\sigma$.

\emph{Finite verification.}
For each of the at most $20^2$ ordered pairs, form
$k=\ell+m$ and accumulate the Gaussian-integer vector
$(\widehat h(\ell)\cdot m)\widehat h(m)$ under the quadratic
amplitude monomial. Eliminate $k=0$ and use
\[
 |\mathbb P_kv|^2
  =|v|^2-\frac{|k\cdot v|^2}{|k|^2}
\]
with rational arithmetic, then combine monomials and conjugate
terms. The same computation tests sector membership and the
$|k|^2$ partition. The stand-alone standard-library verifier
\texttt{tools/verify\_attack24\_exact.py} carries out all 20-mode
and 108-mode checks with integer Gaussian coefficients and
\texttt{Fraction} denominators. A separate physical-space
pseudospectral implementation reproduces
\eqref{eq:AT-leakage} to floating-point round-off for three
independent amplitude choices. The finite enumeration is an exact
computer-assisted identity, not an analytic theorem concerning
infinite spectral evolution.

\paragraph{Initial energy transfer and viscosity.}
For the isolated unforced Navier--Stokes evolution with initial
datum $h$, local smooth existence gives
$\widehat u(k,t_0)=0$ on $\mathcal N$ and
\[
 \partial_t\widehat u(k,t_0)
   =-2\pi i\,\mathbb P_k\mathcal B_k(h,h)
   \quad(k\in\mathcal N).
\]
Consequently, for $s\downarrow0$,
\begin{equation}\label{eq:AT-shorttime}
 \sum_{k\in\mathcal N}|\widehat u(k,t_0+s)|^2
 =4\pi^2s^2\mathcal L(A,B,C)+o(s^2).
\end{equation}
This is nonzero initial leakage, not a monotone shell-to-shell
cascade or a lower bound at terminal times.

Parseval gives exactly
\begin{equation}\label{eq:AT-energy}
 \|h\|_2^2=4A^2+40B^2+40C^2,\qquad
 \|\nabla h\|_2^2=(2\pi)^2(4A^2+200B^2+240C^2).
\end{equation}
The instantaneous nonlinear energy derivatives in the three
occupied wavevector families $|k|^2=1,5,6$, respectively, are
\[
 0,\qquad -320\pi\sigma ABC,\qquad
 +320\pi\sigma ABC.
\]
Their sum is zero, as required by incompressibility.
In particular, for each of the eight receiver modes of type $r$,
\[
 \frac{d}{dt}|\widehat u(r)|^2\Big|_{t=t_0}
   =40\pi\sigma ABC-240\nu\pi^2 C^2.
\]
For $A,B,C>0$, gain against viscosity at this occupied receiver
requires $\sigma=+1$ and $AB>6\pi\nu C$. The choice
$\sigma=-1$ used in the original first-mediator calculation
has the opposite occupied-receiver energy-transfer direction,
although it yields the same off-support polynomial
\eqref{eq:AT-leakage}. This phase distinction is essential.

\noindent\textbf{Proposition AT.2 (second-generation exit).}
Let $k_*=(3,2,1)$ and assume the isolated unforced
Navier--Stokes evolution begins with the above seed. Then
$k_*\notin\mathcal S\cup\mathcal N$, but
\begin{equation}\label{eq:AT-second}
 \partial_t^2\widehat u(k_*,t_0)
  =16\pi^2\sigma A^2C\,b .
\end{equation}
In particular, if $AC\ne0$, the initial 20 modes together
with all 108 first-generation modes still fail to form a
Fourier-support invariant subsystem.

\emph{Proof.}
Write $s=p+r=(2,2,1)$. The first-derivative identity is
$\partial_t\widehat u(s,t_0)=4\pi\sigma AC\,b$.
At $k_*=p+s$ the only contributing ordered
pair involving an initially generated mode is $(p,s)$ and
its reverse. The reverse derivative term vanishes because
$b\cdot p=0$, whereas
\[
 (\widehat h(p)\cdot s)\,\partial_t\widehat u(s,t_0)
 =(2iA)(4\pi\sigma AC\,b)
 =8\pi i\sigma A^2C\,b .
\]
Since $k_*\cdot b=0$, the Leray projection leaves
this vector unchanged. The outer multiplier $-2\pi i$
gives \eqref{eq:AT-second}. At $k_*$ both the initial
amplitude and first derivative vanish, so the viscous
term contributes nothing to the second derivative.
Consequently
$\widehat u(k_*,t_0+s)
=8\pi^2\sigma A^2C\,b\,s^2+o(s^2)$.
This conclusion concerns the isolated initial evolution.
An independent physical-space pseudospectral
second-derivative calculation reproduces
\eqref{eq:AT-second} at the tested amplitudes.

\paragraph{A carrier-inclusive necessary test.}
Let $U$ be any smooth periodic divergence-free carrier at
$t_0<T$ and write $u=U+w$. Define its exact projected residual
\[
 \mathcal R_\nu(U)
 =\partial_tU-\nu\Delta U+
   \mathbb P\nabla\cdot(U\otimes U)
\]
and the carrier-coupling source
\[
 \mathcal K_U(w)
 =-\mathcal R_\nu(U)
  -\mathbb P\nabla\cdot(U\otimes w+w\otimes U).
\]
If the full $u$ solves the unforced equation, the correction
satisfies
\[
 \partial_tw-\nu\Delta w+
 \mathbb P\nabla\cdot(w\otimes w)=\mathcal K_U(w).
\]
Suppose $w(t_0)=h$, so its coefficients vanish on $\mathcal N$.
Writing $\Pi_{\mathcal N}$ for the orthogonal Fourier projector,
the triangle inequality and \eqref{eq:AT-leakage} imply
\begin{equation}\label{eq:AT-carrier}
 \left\|\Pi_{\mathcal N}\partial_tw(t_0)\right\|_2
 \ge\left[
 2\pi\sqrt{\mathcal L(A,B,C)}
 -\left\|\Pi_{\mathcal N}\mathcal K_U(h)(t_0)\right\|_2
 \right]_+ .
\end{equation}
Here $[x]_+=\max(x,0)$.
An explicit, generally non-sharp bound is
\[
 \|\Pi_{\mathcal N}\mathcal K_U(h)\|_2
 \le \|\mathcal R_\nu(U)\|_2
    +\|U\|_\infty\|\nabla h\|_2
    +\|\nabla U\|_\infty\|h\|_2 .
\]
Thus first-generation leakage persists in the full correction
whenever the actual projected carrier coupling is strictly
smaller than $2\pi\sqrt{\mathcal L}$; cancellation is not
excluded when it is comparable or larger. No value for the
actual frozen Version~17 carrier-coupling norm is inferred
from this finite computation.

Finally, the exact correction energy relation is
\begin{equation}\label{eq:AT-wenergy}
 \frac12\frac{d}{dt}\|w\|_2^2+\nu\|\nabla w\|_2^2
 =-\int_{\mathbb T^3}w\cdot(w\cdot\nabla U)\,dx
   -\langle\mathcal R_\nu(U),w\rangle ,
\end{equation}
whereas the full unforced solution must obey
$\frac12\frac{d}{dt}\|U+w\|_2^2+
\nu\|\nabla(U+w)\|_2^2=0$.
The finite-mode identities do not control
\eqref{eq:AT-wenergy} over a shrinking terminal interval.

\paragraph{Scope of the result.}
The exact data show that distinct fourfold sectors generate
many new modes, including both higher and lower frequencies,
and they quantify the initial transfer. They do not supply
a sustained nonlinear cascade, an energy-compatible
terminal-time modulation, a bound on
$\Pi_{\mathcal N}\mathcal K_U$, or non-cancellation of the
frozen singular trace. The smooth unrotated invariant
shear-advection subsystem from the preceding finite-triad
calculation also remains an obstruction to using a single
polarization ray as a blow-up mechanism. In particular, no
unforced finite-time singularity theorem follows from
Appendix~AT.
